\section{Rubrik Penilaian Komprehensif}

Penilaian komprehensif memerlukan pendekatan yang sistematis dan objektif. Menurut sumber evaluasi, rubrik yang baik harus mencakup berbagai aspek:
\begin{itemize}
  \item \textbf{Kriteria penilaian}: Kriteria harus jelas, terukur, dan dapat diukur secara konsisten.
  \item \textbf{Feedback konstruktif}: Memberikan umpan balik yang mendalam dan konstruktif untuk perbaikan.
  \item \textbf{Authenticity}: Memastikan keaslian penilaian dan mencegah plagiarisme.
  \item \textbf{Fairness}: Menjamin proses penilaian yang adil dan transparan.
\end{itemize}

Rubrik penilaian yang baik membantu meningkatkan kualitas pembelajaran dan memastikan pencapaian yang objektif dan terukur. Menurut sumber evaluasi, implementasi rubrik yang konsisten dan berbasis bukti sangat penting untuk pengembangan program studi yang berkelanjutan. Pendekatan yang sistematis dan objektif memastikan bahwa evaluasi tidak hanya mengukur hasil belajar, tetapi juga mengembangkan kemampuan profesional dan etika akademis mahasiswa.

\subsection{Metode Penilaian}

Berbagai metode penilaian dapat digunakan untuk mengukur pencapaian:
\begin{itemize}
  \item \textbf{Observasi}: Pengamatan langsung proses pembelajaran dan partisipasi mahasiswa.
  \item \textbf{Portofolio}: Evaluasi karya mahasiswa untuk menunjukkan kemajuan praktikum.
  \item \textbf{Presentasi}: Presentasi hasil proyek atau penelitian.
  \item \textbf{Tes Tertulis}: Uji pemahaman dengan soal esai atau praktikum.
  \item \textbf{Peer Review}: Evaluasi kerja mahasiswa oleh teman sekelas.
  \item \textbf{Self-Assessment}: Refleksi diri dari pembelajaran.
\end{itemize}

Metode penilaian yang baik dirancang untuk memberikan feedback yang konstruktif dan mendukung pengembangan. Menurut sumber evaluasi, kombinasi berbagai metode ini memberikan gambaran yang komprehensif tentang kinerja dan pemahaman mahasiswa, serta memfasilitasi perbaikan berkelanjutan.

\textbf{Kriteria nilai}: A (85--100), B (70--84), C (60--69), D (50--59), E (<50).
